% paper/sections/related_work.tex

\section{Related Work}
\label{sec:related}

\subsection{Domain Randomisation for Sim-to-Real Transfer}

Domain randomisation (\DR{}) trains policies across distributions of
simulated environments, hypothesising that the real world falls within
the training distribution~\cite{tobin2017}.
\citet{tobin2017} first demonstrated visual \DR{} for robotic
grasping, achieving zero-shot sim-to-real transfer from randomised
rendered images.
\citet{peng2018} extended \DR{} to dynamics---mass, friction, joint
parameters---for legged locomotion, establishing randomised physics as a
viable route to robust transfer.

Adaptive approaches that update randomisation ranges during training
have shown further gains.
OpenAI's \textsc{Adr}~\cite{akkaya2019} expands parameter boundaries
when per-parameter performance buffers indicate competence, enabling
dexterous manipulation with 132 randomised parameters.
\textsc{Doraemon}~\cite{tiboni2024} frames adaptive \DR{} as entropy
maximisation subject to a success constraint, using Bayesian
optimisation to select parameter distributions.

Our \CDR{} adapts the \textsc{Adr} concept to underwater fluid physics,
expanding all six parameters simultaneously with a single shared
success window rather than per-parameter buffers.
Unlike \textsc{Adr} and \textsc{Doraemon}---which target manipulation or
locomotion on ground-based systems---we apply curriculum \DR{} to the
qualitatively distinct parameter space of hydrodynamic forces and
evaluate energy efficiency as a primary metric; both are absent from
prior \DR{} literature.

\subsection{Reinforcement Learning for AUV Control}

Deep \RL{} has been applied to \AUV{} station-keeping~\cite{carlucho2018},
depth and heading control~\cite{shi2020}, and three-dimensional
navigation.
\citet{arndt2024} demonstrated zero-shot sim-to-real transfer for a
6-DOF \AUV{} using \SAC{} with Uniform \DR{} on two physical parameters
(centre of buoyancy and displaced volume), achieving performance
comparable to hand-tuned \PID{} on a real vehicle.
\citet{chaffre2025} combined \SAC{} with bio-inspired experience replay
and enhanced Uniform \DR{} for \AUV{} station-keeping under current
disturbances, with physical tank validation.

\textsc{MarineGym}~\cite{chu2025} is the most comprehensive underwater
\RL{} benchmark to date: GPU-accelerated simulation (NVIDIA Isaac Sim)
with five \AUV{} models, three task types, and a Uniform \DR{} toolkit.
However, \textsc{MarineGym} studies neither curriculum scheduling of
\DR{} nor energy efficiency, and uses proprietary simulation rather than
the open-source \MuJoCo{} we employ.
No prior work quantifies \PID{} fragility under fluid-parameter
distribution shift---a $97$-pp collapse we demonstrate here for the
first time.

\subsection{Curriculum Learning}

Curriculum learning~\cite{bengio2009}---training on examples ordered
from easy to hard---improves generalisation and convergence in
supervised and reinforcement learning.
Self-paced \RL{}~\cite{klink2020} adapts task difficulty to the
learner's current capability using contextual MDPs.
Our work is, to our knowledge, the first to apply performance-gated
curriculum \DR{} to underwater fluid physics and the first to
demonstrate that curriculum scheduling reduces energy consumption
relative to Uniform \DR{}---a result with direct operational significance
for battery-constrained \AUV{} deployments.